\scanpage{5}
Descriptions conjecturales des motifs sur corps de nbs \deleted{\unsure{$F$|$K$}} $k$

1) Double réseau Betti-De Rham

2) Double réseau Betti-Hodge

3) \deleted{Betti + \unsure{bigradé|bigrad.}} structure de Hodge $(k=\overline{\unsure{Q|\mathbb Q}})$

4) \deleted{Betti + \unsure{connexion} réseau Betti} Betti-(Gauss)Manin $(k=\overline Q)$

5) \deleted{Tate} Betti-Tate

6) \deleted{De Rham - cristal} De Rham filtré cristallin

Descriptions conjecturales des groupes $p$-divisibles \inserted{mod isogénie} (plus \unsure{génér\textsuperscript{t}|génér\textsuperscript{alt}|gén.} des pseudo-motifs) sur un corps local [\deleted{\unsure{spectre}} \inserted{corps de fr.} d'un ann. de val. discrète complet à corps rés \deleted{\unsure{parfait}}] :

$3'$) structure de Hodge-\inserted{Tate\unsure{;|:} module $N$ de t.f. sur $Q_p$, et bigrad. sur $N\otimes_{\unsure{Q_p|Q|\mathbb Q_p}} C$} (mod ext. finie de $K$)

$2'$) double réseau Tate-Hodge

$5'$) Tate (module galoisien $\unsure{l|\ell}$-adique)

$6'$) \deleted{\unsure{module de vect. de Di|module vect. de Di|module Vect. de di}} cristal filtré [vectoriel \deleted{\unsure{sur}} dim finie sur $K$, une application $\sigma$-semi-linéaire, \deleted{\unsure{et|+}} filtration, graduation]

Par quoi remplacer $1',4'$ ??
% archive sheet number: 4
